\section{Metaphysical Knowledge}

The following quotes are taken from \textit{Introduction to the Study of the Hindu Doctrines} by \textbf{Rene Guenon}. They characterize the nature of the metaphysical enterprise.

\paragraph{Always and Everywhere}


The object of metaphysical knowledge must always be absolutely the same and can in no way be something that changes or that is subject to the influences of time and place. That is, it is true always and everywhere. 

\paragraph{Completeness}


It is absolutely impossible to make any `discoveries' in metaphysics, for in a type of knowledge which calls for the use of no specialized or external means of investigation, all that is capable of being known may have been known by certain persons at any and every period; and this in fact emerges clearly from a profound study of the traditional metaphysical doctrines.

\paragraph{Immutability}


The metaphysical point of view is radically opposed to the historical point of view, or what passes for such, and this opposition will be seen to amount not only to a question of method, but also, what is far more important, to a real question of principle, since the metaphysical point of view, in its essential immutability, is the very negation of the notions of evolution and progress.

\paragraph{Certitude}


No notice must be taken of contingencies such as individual influences, which are strictly non-existent from this point of view and cannot affect the doctrine in any way; the latter, being of the universal order, is thereby essentially supra-individual, and necessarily remains untouched by such influences. Even circumstances of time and space, we must repeat, can only affect the outward expression but not the essence of the doctrine; moreover there can be no question here, as there is in the relative and contingent order, of beliefs or opinions; metaphysical knowledge essentially implies permanent and changeless certitude.

\paragraph{Uncontestable}


Metaphysics of necessity excludes every conception of a hypothetical character, whence it follows that metaphysical truths, in themselves, cannot in any way be contestable. Consequently, if there sometimes is occasion for discussion and controversy, this only happens as a result of a defect in exposition or of an imperfect comprehension of those truths.

\paragraph{Ineffability}


Every exposition possible of metaphysics is necessarily defective, because metaphysical conceptions, by reason of their universality, can never be completely expressed, nor even imagined, since their essence is attainable by the pure and formless intelligence alone; they vastly exceed all possible forms, especially the formulas in which language tries to enclose them, which are always inadequate and tend to restrict their scope and therefore to distort them.

\paragraph{Knowledge of the Universal}


Knowledge belonging to the universal order of necessity lies beyond all the distinctions that condition the knowledge of individual things. The object of metaphysics is in no way comparable to the particular object of any other kind of knowledge whatsoever.

\paragraph{Being and knowledge are unified}


When one speaks of the means of attaining metaphysical knowledge, it is evident that such means can only be one and the same thing as knowledge itself, in which subject and object are essentially unified. That is, metaphysical knowledge is possible only with a change in one's level of being or state of consciousness. 

\paragraph{Non-discursive}


The means of achieving metaphysical knowledge cannot in any way resemble the exercise of a discursive faculty such as individual human reason.

\paragraph{Beyond Rational}


We are dealing with the supra-individual and consequently with the supra rational order, which does not in any way mean the irrational: metaphysics cannot contradict reason, but it stands above reason, which is only a secondary means for the formulation and external expression of truths that lie beyond its province and outside its scope.

\paragraph{Intuition}


Metaphysical truths can only be conceived by the use of a faculty that does not belong to the individual order, and that, because of the immediate character of its operation, may be called \emph{intuitive}.

\paragraph{Direct knowledge of principles}


The intellect is that faculty which possesses a direct knowledge of principles.

\paragraph{Truer than science}


Aristotle declares that ``the intellect is truer than science'', which amounts to saying that it is truer than the reason which constructs that science.

\paragraph{Infallible}


It is necessarily infallible because of the fact that its operation is immediate and, not being really distinct from its object, it is identified with the truth itself.

These characteristics are perfectly clear to someone who understands them, but are opaque to someone who does not. But to apply them to particular situations: neither the individual, nor mankind as a whole, ``evolves'' into higher knowledge. Debates about particular exterior forms are worse than futile. It is transcendent and hence not reducible to any biological process.

In conclusion, there is no point in ``debating'' any of these characteristics.


\flrightit{Posted on 2010-05-10 by Cologero}

\begin{center}* * *\end{center}

\begin{footnotesize}
\begin{sffamily}
\texttt{Sati Shankar on 2010-05-11 at 01:06 said:}

The riddle creeps in as soon as we start expressing, with words or signs, the inexpressible.What Guenon says of the ultimate metaphysical reality, has been expresses only as neti, neti,... (not this .not this... ..). The attributes assigned above to that are in themselves,dualistic limitations imposed. Once we take the point of ``That Unity'' almost all qualifications and discriminating signs are vanished.Then there is no phenomenological attribute attached, then all what we call,knowledge, hierarchy,evolution, time, space,even consciousness and hence ``intuition'' of living beings cease to be.

With this background arises a question which becomes self limiting phenomenologically.The ``tradition'' which Guenon lead,and specially the ``traditionists'' who follow and pursue, as distinct from''traditionalists'' to use the Guenon's terminology, who plea to extract the tradition, being expressions of revelations and wish to assign, and to purge out the deviations from primordial tradition as corruptions in expression of those revelations, based on''Rational'' grounds.

This ``Rational Grounds'' is based on ``Reason'',(ratio),which itself presumes duality. It needs ordering in time and space and hence is evolutionary.

The question is ``can'' That which can be expressed NOT other than by neti... neti... the method of negation, be judged, edited and filtered by the ``Rational'' methods to apply to the world Traditions?The Tradition which is really timeless, akAliko, eternal, sAshvat is that which embraces and pervades ``all'' is the Primeordial Tradition of ``Unity''.

Any phenomenological editing,of that metaphysical, \{which is impossible in true sense when viewed from a dualistic standpoint) is bound to allow creeping the mental arbitrations, adulterations in the name of intuitions and the rational.It would emerge as a new sect or religion with the same possibilities of use and abuse as with other Religions. The question is who authorizes, in a Primordial Tradition, to do it?

\hfill

\texttt{Cologero on 2010-05-11 at 08:37 said:}

Guenon is writing for those who may be lacking any Traditional knowledge. Therefore, he is trying to make them aware that it is something that surpasses any rational discussion of it. After all, ``ineffability'' is indeed in the list. There is no intent, either on Guenon's part or mine, to reduce the metaphysical to the rational. We have written about the degrees of knowledge several times.

Also, there is no attempt to create a Traditionalism apart from the Traditions. That is addressed
\textit{here}\footnote{\url{http://www.gornahoor.net/?p=383}}.


I disagree with this: ``Then there is no phenomenological attribute attached, then all what we call,knowledge, hierarchy,evolution, time, space,even consciousness and hence intuition of living beings cease to be.'' If they never were, how could they cease to be?

That is because ``phenomenological attributes'' appear to be. There is no ``taking the point of Unity'', as that implies that at one point in time there was no unity, and then there now is unity. Not at all ... what appears to be, appears to be. What is, is.

\hfill

\end{sffamily}
\end{footnotesize}

